\documentclass[11pt,a4paper]{article}

\usepackage[T1]{fontenc}
\usepackage[spanish,es-nodecimaldot]{babel}
\usepackage{lmodern}
\usepackage{amsmath,amssymb,amsthm,mathtools,bm}
\usepackage{booktabs,array}
\usepackage{microtype}
\usepackage[margin=2.35cm]{geometry}
\usepackage{xcolor}
\usepackage{hyperref}
\hypersetup{colorlinks=true,linkcolor=blue!45!black,citecolor=blue!45!black,urlcolor=blue!55!black}

\newcommand{\Tr}{\operatorname{Tr}}
\newcommand{\id}{\mathbb I}
\newcommand{\Tfun}{\mathfrak T}
\newcommand{\W}{W}
\newcommand{\Fstar}{F_{\star}}

\newcommand{\Ngstar}{\mathcal N_{\mathrm{GME}}^{\star}}
\newcommand{\Pstar}{\mathcal P_{\star}}
\newcommand{\starprod}{\star}
\newcommand{\dd}{\mathrm d}
\newcommand{\e}{\mathrm e}

\newtheorem{theorem}{Teorema}
\newtheorem{proposition}[theorem]{Proposición}
\newtheorem{lemma}[theorem]{Lema}
\theoremstyle{definition}
\newtheorem{definition}[theorem]{Definición}
\newtheorem{remark}[theorem]{Observación}

\title{Detector jerárquico de entrelazamiento multiqubit\\
desde funciones de Wigner Stratonovich--Weyl}
\author{Enrique E. Casanova Benítez\\
\small Universidad Autónoma de Santo Domingo (UASD), República Dominicana
\and
Oscar Jasel Berra Montiel\\
\small Universidad Autónoma de San Luis Potosí (UASLP), México}
\date{\today}

\begin{document}
\maketitle

\begin{abstract}
Construimos un detector jerárquico de entrelazamiento multiqubit cuyo dato de
entrada es la función de Wigner de Stratonovich--Weyl (SW), no una matriz de
densidad. La primera capa refleja las esferas de Bloch asociadas a una
bipartición, calcula momentos mediante el producto estrella finito y reconstruye
el espectro de la transposición parcial con identidades de Newton. El funcional
$F_{\star}^{(S)}$ resultante es un criterio NPT exacto para cada corte, pero no
certifica por sí solo entrelazamiento genuino multipartito (GME) de estados
mixtos. Para cerrar esa brecha introducimos dos capas globales. La primera usa
momentos estrella de mapas positivos sobre todo estado biseparable y define
$\mathcal G_{\star,\ell}$ como la parte negativa del menor autovalor de una
matriz de Hankel. Esta condición matricial es más informativa que utilizar solo
el determinante de Hankel. La segunda formula enteramente en el álgebra SW el
programa semidefinido de testigos completamente descomponibles y produce
$\mathcal N_{\mathrm{GME}}^{\star}$, equivalente a la negatividad genuina y
positivo exactamente fuera del conjunto de mezclas PPT. Verificamos la cadena
desde símbolos de Wigner explícitos para GHZ$_3$, W$_3$, GHZ$_4$, familias con
ruido blanco y una mezcla biseparable que es NPT en los tres cortes. GHZ$_4$
se evalúa sobre sus siete biparticiones no redundantes. El error máximo entre los
cálculos estrella y comprobaciones matriciales independientes es
$9.1\times10^{-13}$. El resultado no se presenta como una nueva medida universal:
$F_{\star}$ diagnostica NPT por corte, $\mathcal G_{\star,\ell}$ es un testigo
suficiente dependiente del mapa y $\mathcal N_{\mathrm{GME}}^{\star}$ hereda las
fortalezas y limitaciones de la relajación por mezclas PPT.
\end{abstract}

\noindent\textbf{Palabras clave:} función de Wigner; producto estrella;
transposición parcial; entrelazamiento genuino multipartito; mezclas PPT;
programación semidefinida.

\section{Introducción}

La representación de Wigner reemplaza operadores por funciones sobre un espacio
de fase, pero conserva la no conmutatividad mediante un producto estrella
\cite{Wigner1932,Groenewold1946,Moyal1949}. Para espines y sistemas de dimensión
finita, la correspondencia de Stratonovich--Weyl proporciona un kernel
autodual y covariante \cite{Stratonovich1957,Tilma2016ArbitraryWigner,
RundleEveritt2021Overview}. Esta formulación se ha usado para dinámica,
protocolos cuánticos y validación de correlaciones; en particular, el programa
de cuantización por deformación de Berra-Montiel y colaboradores desarrolla el
producto estrella de qubits y sus exponenciales
\cite{SanchezCordovaBerraMontiel2024QET,BerraMontiel2026StarDynamics}.

La negatividad puntual o el volumen negativo de una función de Wigner puede ser
un testigo útil en familias concretas \cite{Arkhipov2018NegativityVolume}, pero no
es, en general, una medida de entrelazamiento. Un estado separable puede tener
una función cuasiprobabilística negativa y la no negatividad puntual tampoco
garantiza separabilidad. Por ello buscamos funcionales que dependan de la
estructura no local codificada en $W$ y del producto estrella, no solo del signo
ordinario de $W$.

Para dos qubits, la transposición parcial (PT) decide separabilidad
\cite{Peres1996Separability,Horodecki1996Separability}. En el espacio de fase SW,
PT es una reflexión de la componente $y$ de la esfera local y sus momentos se
obtienen con potencias estrella
\cite{Elben2020PTMoments,YuGuehne2021Hankel,Zhang2022PTMoments}. En sistemas
mayores aparecen dos obstáculos.
Primero, el determinante de la PT pierde información: varios autovalores
negativos pueden dar determinante positivo y la deficiencia de rango puede
forzar determinante cero. Segundo, aun una detección NPT exacta en todos los
cortes no implica GME para estados mixtos. La razón es convexa: una mezcla de
estados separables en cortes diferentes puede ser NPT en cada corte y seguir
siendo biseparable \cite{GuhneToth2009Review,Jungnitsch2011PPTMixtures}.

En este trabajo proponemos una arquitectura de tres capas:
\begin{enumerate}
\item $F_{\star}^{(S)}[W]$: reconstrucción espectral exacta y diagnóstico NPT en
cada bipartición $S|S^c$;
\item $\mathcal G_{\star,\ell}^{(\mathcal M)}[W]$: certificado GME de bajo orden
obtenido de la positividad completa de matrices de Hankel asociadas a un mapa
GME;
\item $\mathcal N_{\mathrm{GME}}^{\star}[W]$: programa semidefinido SW de
testigos completamente descomponibles, usado como capa global con información
completa.
\end{enumerate}
La segunda capa traduce a espacio de fase los mapas GME y sus momentos
\cite{HuberSengupta2014PositiveMaps,Clivaz2017GMEMaps,
Mukherjee2025GMEMoments}. Nuestra elección de la parte negativa del menor
autovalor de Hankel usa la condición demostrada, $H_\ell\succeq0$, en lugar de
reducirla al escalar $\det H_\ell\geq0$. La tercera capa traduce la relajación de
mezclas PPT de Jungnitsch, Moroder y Gühne \cite{Jungnitsch2011PPTMixtures} al
cono positivo del álgebra estrella.

La contribución es una síntesis SW verificable y una lógica de decisión que
evita dos sobrerreclamos comunes: ni ``NPT en todos los cortes'' significa GME
mixta, ni pertenecer al conjunto de mezclas PPT excluye todos los estados GME.
La formulación es exacta como cambio de representación, pero no reduce por sí
sola el costo de adquirir una función $W$ informacionalmente completa.

\section{Álgebra de Wigner para \texorpdfstring{$N$}{N} qubits}

\subsection{Kernel, medida y funcional de traza}

Fijamos para cada qubit
\begin{equation}
 \Delta(\bm n)=\frac12\left(\id+\sqrt3\,\bm n\cdot\bm\sigma\right),
 \qquad \bm n\in S^2,
 \label{eq:kernel}
\end{equation}
con
\begin{equation}
 \dd\mu(\bm n)=\frac{\sin\theta\,\dd\theta\,\dd\phi}{4\pi},
 \qquad \int_{S^2}\dd\mu=1,
 \qquad \int_{S^2}\Delta\,\dd\mu=\frac{\id}{2}.
 \label{eq:measure}
\end{equation}
Para un operador $A$ sobre $N$ qubits definimos
\begin{equation}
 W_A(\bm n_1,\ldots,\bm n_N)
 =\Tr\!\left[A\,\Delta(\bm n_1)\otimes\cdots\otimes\Delta(\bm n_N)\right].
 \label{eq:SWmap}
\end{equation}
La inversa y el funcional de traza son
\begin{align}
 A&=2^N\int W_A\,\Delta(\bm n_1)\otimes\cdots\otimes\Delta(\bm n_N)
       \prod_{r=1}^{N}\dd\mu_r,\label{eq:inverse}\\
 \Tfun_N[f]&:=2^N\int f\prod_{r=1}^{N}\dd\mu_r,
 \qquad \Tfun_N[W_A]=\Tr A.\label{eq:tracefun}
\end{align}
Los factores $2^N$ no son opcionales: compensan exactamente
$\int\Delta\dd\mu=\id/2$.

\subsection{Producto estrella finito}

El producto estrella se define por
\begin{equation}
 W_A\starprod W_B:=W_{AB}.
 \label{eq:star-def}
\end{equation}
En la base local
\begin{equation}
 e_0=1,\qquad e_i=\sqrt3\,n_i=W_{\sigma_i},
 \label{eq:e-basis}
\end{equation}
su tabla es finita:
\begin{equation}
 e_0\starprod e_\mu=e_\mu\starprod e_0=e_\mu,
 \qquad
 e_i\starprod e_j=\delta_{ij}e_0+i\epsilon_{ijk}e_k.
 \label{eq:star-table}
\end{equation}
La regla se factoriza sobre los qubits. Si
$\bm\mu=(\mu_1,\ldots,\mu_N)$ y
$e_{\bm\mu}=\prod_r e_{\mu_r}(\bm n_r)$, entonces
\begin{equation}
 W_A=\sum_{\bm\mu}c_{\bm\mu}[A]e_{\bm\mu},
 \qquad
 c_{\bm\mu}[A]=2^{-N}\Tr(A\sigma_{\bm\mu}),
 \qquad
 A=\sum_{\bm\mu}c_{\bm\mu}[A]\sigma_{\bm\mu}.
 \label{eq:coefficients}
\end{equation}
En particular,
\begin{equation}
 \Tfun_N[W_A\starprod W_B]=\Tr(AB).
 \label{eq:star-trace}
\end{equation}

\subsection{La entrada operacional es \texorpdfstring{$W$}{W}}

El algoritmo recibe los coeficientes $c_{\bm\mu}[W]$ o una representación
equivalente de $W$. No requiere que se proporcione una matriz $\rho$. Para
formular positividad sin abandonar el espacio de fase introducimos el cono
\begin{equation}
 \Pstar:=\left\{f=f^*:\;
 \Tfun_N[g^*\starprod f\starprod g]\geq0
 \ \text{para todo }g\right\}.
 \label{eq:star-positive-cone}
\end{equation}
Escribimos $f\succeq_\star0$ si $f\in\Pstar$, y
$0\preceq_\star f\preceq_\star1$ si $f$ y $1-f$ pertenecen a $\Pstar$.
Por la biyección SW, $f\succeq_\star0$ es equivalente a positividad del elemento
del álgebra representado por sus coeficientes. Así, una entrada física satisface
\begin{equation}
 W=W^*,\qquad W\succeq_\star0,\qquad \Tfun_N[W]=1.
 \label{eq:physical-W}
\end{equation}
Esta positividad estrella no es positividad puntual: $W(\Omega)$ puede tomar
valores negativos.

\section{Capa I: espectroscopia NPT por corte}

\subsection{Transposición parcial como reflexión}

En la base computacional, $\sigma_x^T=\sigma_x$,
$\sigma_y^T=-\sigma_y$ y $\sigma_z^T=\sigma_z$. Para una bipartición
$S|S^c$, definimos
\begin{equation}
 (\Theta_SW)(\bm n_1,\ldots,\bm n_N)
 =W(\bm n'_1,\ldots,\bm n'_N),
 \qquad
 \bm n'_r=\begin{cases}(n_{rx},-n_{ry},n_{rz}),&r\in S,\\
 \bm n_r,&r\notin S.
 \end{cases}
 \label{eq:reflection}
\end{equation}
Entonces $\Theta_SW_A=W_{A^{T_S}}$. En coeficientes basta cambiar el signo de
cada término que contenga un número impar de índices $y$ dentro de $S$.

\subsection{Momentos estrella e identidades de Newton}

Sea $d=2^N$. Los momentos reflejados son
\begin{equation}
 q_n^{(S)}[W]
 :=\Tfun_N\!\left[(\Theta_SW)^{\starprod n}\right],
 \qquad n=1,\ldots,d.
 \label{eq:q-moments}
\end{equation}
Si $\lambda_1^{(S)},\ldots,\lambda_d^{(S)}$ son los valores espectrales del
elemento reflejado, entonces
\begin{equation}
 q_n^{(S)}=\sum_{j=1}^{d}(\lambda_j^{(S)})^n.
 \label{eq:power-sums}
\end{equation}
Los polinomios simétricos elementales se reconstruyen por
\begin{equation}
 e_0^{(S)}=1,\qquad
 k e_k^{(S)}=\sum_{j=1}^{k}(-1)^{j-1}e_{k-j}^{(S)}q_j^{(S)},
 \qquad k=1,\ldots,d.
 \label{eq:newton}
\end{equation}

\begin{theorem}[Criterio NPT estrella exacto]
Para cualquier entrada física $W$ y cualquier corte $S|S^c$, defínase
\begin{equation}
 \Fstar^{(S)}[W]:=\min_{1\leq k\leq d}e_k^{(S)}[W].
 \label{eq:Fstar}
\end{equation}
Entonces
\begin{equation}
 \Fstar^{(S)}[W]\geq0
 \iff \Theta_SW\succeq_\star0,
 \qquad
 \Fstar^{(S)}[W]<0
 \iff \text{$W$ es NPT en el corte $S|S^c$}.
 \label{eq:Fstar-theorem}
\end{equation}
\end{theorem}

\begin{proof}
La reflexión produce un elemento hermítico, por lo que sus raíces espectrales
son reales. Si todas son no negativas, todos los $e_k$ son no negativos.
Recíprocamente,
\begin{equation}
 P(t)=\prod_{j=1}^{d}(t+\lambda_j)
 =t^d+e_1t^{d-1}+\cdots+e_d.
\end{equation}
Si todos los coeficientes son no negativos, $P(t)>0$ para todo $t>0$.
Una raíz espectral negativa produciría la raíz positiva
$t=-\lambda_j$, contradicción. Por tanto, todos los $e_k\geq0$ si y solo si el
elemento reflejado es positivo.
\end{proof}

Para dos qubits, $e_4=\det(\rho^{T_B})$ tiene signo negativo exactamente en el
caso entrelazado. Esta simplificación no sobrevive a dimensiones mayores. En un
corte $2\otimes4$ puede haber hasta tres autovalores PT negativos
\cite{Rana2013PartialTranspose,Johnston2013PTBound}; un número par de ellos da
determinante positivo, y un cero da determinante nulo. Para GHZ$_3$, por ejemplo,
el espectro reflejado es
\begin{equation}
 \left\{-\tfrac12,0,0,0,0,\tfrac12,\tfrac12,\tfrac12\right\},
\end{equation}
de modo que el determinante se anula aunque $\Fstar=-1/4$.

\section{Por qué los cortes no resuelven la GME mixta}

Un estado multipartito es biseparable si admite una descomposición convexa
\begin{equation}
 W_{\mathrm{bisep}}=\sum_{S\in\mathfrak B_N}\sum_j
 p_{S,j}\,W_{S,j}^{\mathrm{sep}},
 \label{eq:biseparable}
\end{equation}
donde $\mathfrak B_N$ contiene un representante de cada bipartición y cada
término es separable en su corte, aunque cortes diferentes puedan aparecer en la
misma mezcla. Un estado que no admite esta forma es genuinamente multipartito
entrelazado.

En estados puros, entrelazamiento en todo corte equivale a GME. Para estados
mixtos no: la convexificación en la Ec.~\eqref{eq:biseparable} puede producir una
mezcla NPT en todos los cortes. Además, en dimensiones $2\otimes4$ existen
estados PPT entrelazados \cite{Horodecki1997BoundEntanglement}. Por tanto,
\begin{equation}
 \Fstar^{(S)}<0\ \forall S
 \quad\nRightarrow\quad
 \text{GME para estados mixtos}.
 \label{eq:allcuts-not-gme}
\end{equation}

Una relajación convexa útil son las mezclas PPT:
\begin{equation}
 W_{\mathrm{PPTmix}}=\sum_{S\in\mathfrak B_N}p_S W_S,
 \qquad \Theta_SW_S\succeq_\star0.
 \label{eq:pptmixture}
\end{equation}
Todo estado biseparable es una mezcla PPT, así que demostrar que $W$ no pertenece
a este conjunto certifica GME \cite{Jungnitsch2011PPTMixtures}. El recíproco no
es universal: existen estados GME que pertenecen al conjunto de mezclas PPT;
Ha y Kye construyen uno que es PPT respecto de una bipartición
\cite{HaKye2016PPTGME}. Para estados permutacionalmente invariantes de tres
qubits, en cambio, la relajación caracteriza la biseparabilidad
\cite{Novo2013PermInvariant}.

\section{Capa II: momentos estrella de mapas GME}

\subsection{Mapa de transposición levantado}

Los mapas positivos bipartitos pueden levantarse a mapas hermiticidad-preservantes
que son positivos sobre todo estado biseparable
\cite{HuberSengupta2014PositiveMaps,Clivaz2017GMEMaps}. Para qubits y el mapa de
transposición, su forma SW es
\begin{equation}
 \mathcal M_T^{(N)}[W]
 =\sum_{S\in\mathfrak B_N}\Theta_SW
 +c_N\,\Tfun_N[W],1,
 \qquad
 c_N=\frac{2^{N-1}-2}{2}.
 \label{eq:GME-map}
\end{equation}
La suma usa una sola orientación por bipartición, por lo que
$|\mathfrak B_N|=2^{N-1}-1$. El mapa no preserva la traza: para una entrada
normalizada,
\begin{equation}
 \Tfun_N[\mathcal M_T^{(N)}W]
 =|\mathfrak B_N|+2^Nc_N.
 \label{eq:map-trace}
\end{equation}
Así, el primer momento vale $3+8=11$ para $N=3$ y $7+48=55$ para $N=4$.
Fijar estos valores evita mezclar momentos normalizados y no normalizados.

\begin{lemma}[Positividad sobre el casco biseparable]
Para todo símbolo físico biseparable $W$,
\begin{equation}
 \mathcal M_T^{(N)}[W]\succeq_\star0.
 \label{eq:GME-map-positive}
\end{equation}
La afirmación permanece válida si cada término $\Theta_SW$ es seguido por una
conjugación unitaria local fijada de antemano.
\end{lemma}

\begin{proof}
En cualquier bipartición, todo autovalor de la transposición parcial de un
estado normalizado está acotado inferiormente por $-1/2$. Para un estado puro,
los autovalores potencialmente negativos son $-s_i s_j$, donde $s_i,s_j$ son
coeficientes de Schmidt, y
$s_i s_j\leq(s_i^2+s_j^2)/2\leq1/2$. La cota pasa a mezclas por concavidad del
menor autovalor y no cambia bajo conjugación unitaria.

Considérese primero un componente separable en un corte $S_0|S_0^c$. El término
$\Theta_{S_0}W$ es positivo. Cada uno de los otros
$|\mathfrak B_N|-1=2^{N-1}-2$ términos está acotado por
$-\tfrac12\,1$. En consecuencia,
\begin{equation}
 \sum_{S\in\mathfrak B_N}\Theta_SW
 \succeq_\star-\frac{2^{N-1}-2}{2}\,1,
\end{equation}
y el término $c_N\Tfun_N[W]1$ de la Ec.~\eqref{eq:GME-map} cancela exactamente
esa cota. La linealidad extiende el resultado a todo el casco convexo
biseparable.
\end{proof}

Para adaptar el mapa a ciertas clases puede componerse cada transposición con
una conjugación unitaria local. En el ejemplo GHZ usamos
\begin{equation}
 \widetilde{\mathcal M}_T^{(3)}[W]
 =\sum_{X=A,B,C}\mathcal U_x^{(X)}\Theta_XW+\Tfun_3[W],1,
 \label{eq:modified-map}
\end{equation}
donde $\mathcal U_x$ deja $e_x$ invariante y cambia el signo de $e_y,e_z$ en el
qubit correspondiente. Esta es la representación SW de la transposición seguida
por conjugación con $\sigma_x$ \cite{Clivaz2017GMEMaps,
Mukherjee2025GMEMoments}.

\subsection{Jerarquía de Hankel}

Para un mapa GME admisible $\mathcal M$, definimos
\begin{equation}
 s_n^{(\mathcal M)}[W]
 :=\Tfun_N\!\left[(\mathcal M[W])^{\starprod n}\right],
 \qquad
 [H_\ell^{(\mathcal M)}]_{ij}=s_{i+j+1}^{(\mathcal M)},
 \quad i,j=0,\ldots,\ell.
 \label{eq:GME-moments}
\end{equation}
Se necesitan momentos hasta orden $2\ell+1$.

\begin{theorem}[Certificado GME de Hankel]
Si $W$ es biseparable, entonces
\begin{equation}
 H_\ell^{(\mathcal M)}[W]\succeq0
 \qquad\text{para todo }\ell.
 \label{eq:Hankel-positive}
\end{equation}
En consecuencia,
\begin{equation}
 \mathcal G_{\star,\ell}^{(\mathcal M)}[W]
 :=\max\left\{0,-\lambda_{\min}
 \bigl(H_\ell^{(\mathcal M)}[W]\bigr)\right\}>0
 \quad\Longrightarrow\quad W\text{ es GME}.
 \label{eq:Gstar}
\end{equation}
\end{theorem}

\begin{proof}
Por construcción, $\mathcal M[W]\succeq_\star0$ para toda entrada biseparable.
Sean $\lambda_r\geq0$ sus valores espectrales. Con
$[V_\ell]_{ir}=\lambda_r^i$ se obtiene
\begin{equation}
 H_\ell=V_\ell\,\operatorname{diag}(\lambda_r)\,V_\ell^T\succeq0.
\end{equation}
La contraposición demuestra la segunda afirmación.
\end{proof}

La literatura reciente usa $\det H_\ell<0$ como condición suficiente
\cite{Mukherjee2025GMEMoments}. Sin embargo, la demostración establece la
condición matricial completa $H_\ell\succeq0$. Comprobar
$\lambda_{\min}(H_\ell)<0$ es estrictamente más informativo que comprobar solo
el signo del determinante: un número par de autovalores negativos o un cero
puede ocultar la indefinición. En nuestros datos, $H_3$ del GHZ puro tiene
determinante nulo y menor autovalor negativo. El funcional
\eqref{eq:Gstar} no se reclama como monótono LOCC; es la profundidad de violación
de un testigo de momentos para un mapa y un orden fijados.

\section{Capa III: SDP de mezclas PPT en el álgebra estrella}

\subsection{Testigos completamente descomponibles}

Un símbolo real $w$ es completamente descomponible si, para cada
$S\in\mathfrak B_N$, existen $p_S,q_S$ tales que
\begin{equation}
 w=p_S+\Theta_Sq_S,
 \qquad
 0\preceq_\star p_S\preceq_\star1,
 \qquad
 0\preceq_\star q_S\preceq_\star1.
 \label{eq:fully-decomp}
\end{equation}
Definimos
\begin{align}
 \eta_\star[W]
 &=\min_{w,\{p_S,q_S\}}
 \Tfun_N[w\starprod W]
 \quad\text{sujeto a \eqref{eq:fully-decomp}},\label{eq:eta-sdp}\\
 \Ngstar[W]&:=\max\{0,-\eta_\star[W]\}.
 \label{eq:Ngstar}
\end{align}
Esta normalización traslada al problema genuinamente multipartito la lógica de
la negatividad bipartita basada en la norma de traza
\cite{VidalWerner2002Negativity}, pero optimiza sobre testigos completamente
descomponibles y no sobre un único corte.

\begin{theorem}[Equivalencia con la negatividad genuina]
El programa \eqref{eq:eta-sdp} es la imagen SW del SDP de testigos completamente
descomponibles. Por tanto,
\begin{equation}
 \Ngstar[W]>0
 \iff W\text{ no es una mezcla PPT}
 \Longrightarrow W\text{ es GME}.
 \label{eq:Ngstar-theorem}
\end{equation}
Además, $\Ngstar$ hereda convexidad, invariancia local y monotonía bajo LOCC de
la negatividad genuina normalizada
\cite{Jungnitsch2011PPTMixtures}.
\end{theorem}

\begin{proof}
La regla de traza \eqref{eq:star-trace} transforma el objetivo en el valor
esperado del testigo. La reflexión $\Theta_S$ representa exactamente PT y el
cono $\Pstar$ representa positividad semidefinida. Así, cada restricción
\eqref{eq:fully-decomp} es la imagen uno a uno de
$\widehat w=P_S+Q_S^{T_S}$ con $0\leq P_S,Q_S\leq\id$. Para cualquier componente
PPT en el corte $S$,
\begin{equation}
 \Tfun_N[w\starprod W_S]
 =\Tfun_N[p_S\starprod W_S]
 +\Tfun_N[q_S\starprod\Theta_SW_S]\geq0.
\end{equation}
Por convexidad ocurre lo mismo para toda mezcla PPT. El teorema de separación
convexa proporciona un testigo completamente descomponible para cada punto
fuera de ese conjunto \cite{Jungnitsch2011PPTMixtures}. La normalización acotada
produce el monótono citado.
\end{proof}

En la implementación, si
$f=\sum_{\bm\mu}c_{\bm\mu}e_{\bm\mu}$, la aplicación lineal
\begin{equation}
 \Gamma[f]=\sum_{\bm\mu}c_{\bm\mu}\sigma_{\bm\mu}
 \label{eq:gamma-lmi}
\end{equation}
convierte $f\succeq_\star0$ en la LMI $\Gamma[f]\succeq0$. Esto no cambia el
dato de entrada: el optimizador recibe los coeficientes de $W$ y usa
\eqref{eq:gamma-lmi} como representación finita del cono estrella. El SDP se
resuelve con métodos estándar \cite{VandenbergheBoyd1996SDP}.

\section{Detector jerárquico y lógica de decisión}

Resumimos la salida como
\begin{equation}
 \mathfrak D[W]=\left(
 \{\Fstar^{(S)}\}_{S\in\mathfrak B_N},
 \{\mathcal G_{\star,\ell}^{(\mathcal M)}\}_{\mathcal M,\ell},
 \Ngstar\right).
 \label{eq:hierarchical-detector}
\end{equation}
No comprimimos estas cantidades en un escalar universal, porque responden a
preguntas lógicamente distintas.

\begin{center}
\small
\begin{tabular}{p{0.31\linewidth}p{0.61\linewidth}}
\toprule
Resultado & Conclusión válida \\
\midrule
$\Fstar^{(S)}<0$ & Entrelazamiento NPT certificado en $S|S^c$.\\
$\Fstar^{(S)}\geq0$ & PPT en ese corte; separabilidad no garantizada.\\
$\Fstar^{(S)}<0$ en todo corte, $W$ puro & GME certificado.\\
$\Fstar^{(S)}<0$ en todo corte, $W$ mixto & GME no decidida: puede ser una mezcla biseparable.\\
$\mathcal G_{\star,\ell}^{(\mathcal M)}>0$ & GME certificada por momentos del mapa $\mathcal M$.\\
$\Ngstar>0$ & Estado fuera de mezclas PPT; GME certificada.\\
$\Ngstar=0$ & Mezcla PPT; en general GME todavía inconclusa.\\
\bottomrule
\end{tabular}
\end{center}

El protocolo Wigner-first es:
\begin{enumerate}
\item validar \eqref{eq:physical-W};
\item para cada corte, reflejar $W$, calcular $q_1,\ldots,q_{2^N}$ y evaluar
$\Fstar^{(S)}$;
\item si la entrada es pura, usar los cortes para decidir GME;
\item si es mixta, evaluar una familia especificada de mapas GME y aumentar
$\ell$ hasta detectar o alcanzar el presupuesto de momentos;
\item con información completa, resolver \eqref{eq:eta-sdp};
\item devolver también ``inconcluso'' cuando ninguna implicación anterior sea
aplicable.
\end{enumerate}

La primera capa necesita hasta $2^N$ momentos por corte si se desea completitud
NPT. La segunda puede detenerse en bajo orden, pero su éxito depende del mapa y
del estado. La tercera usa $4^N$ coeficientes y LMIs de tamaño $2^N$, por lo que
no elimina el escalamiento exponencial. Si los momentos se miden directamente
mediante protocolos multicopia o mediciones aleatorias, la capa de Hankel puede
evitar reconstrucción completa \cite{Elben2020PTMoments,YuGuehne2021Hankel,
Neven2021MomentHierarchy,Mukherjee2025GMEMoments}; esa ventaja experimental no se
deduce solo de disponer de una tabla completa de $W$.

\section{Ejemplos desde funciones de Wigner explícitas}

Todos los ejemplos siguientes entran al código como coeficientes de $W$. Las
matrices asociadas se construyen únicamente después, para comprobaciones
independientes y para representar las LMIs del SDP.

\subsection{GHZ de tres qubits}

Con $e_{ir}=\sqrt3 n_{ir}$,
\begin{align}
 W_{\mathrm{GHZ}}=\frac18\bigl[&1+e_{zA}e_{zB}+e_{zA}e_{zC}+e_{zB}e_{zC}
 +e_{xA}e_{xB}e_{xC}\nonumber\\
 &-e_{xA}e_{yB}e_{yC}-e_{yA}e_{xB}e_{yC}
 -e_{yA}e_{yB}e_{xC}\bigr].
 \label{eq:W-GHZ}
\end{align}
La capa de cortes da $\Fstar^{(A)}=\Fstar^{(B)}=\Fstar^{(C)}=-1/4$.
El mapa modificado \eqref{eq:modified-map} tiene menor autovalor $-1/2$ y
$H_1$ ya es indefinida, con $\lambda_{\min}(H_1)=-0.1439418$. El SDP produce
$\Ngstar=0.5$.

\subsection{Estado W de tres qubits}

Una forma compacta es
\begin{align}
 W_{\mathrm W}=\frac18\bigg[&1+\frac13\sum_i e_{zi}
 -\frac13\sum_{i<j}e_{zi}e_{zj}-e_{zA}e_{zB}e_{zC}\nonumber\\
 &+\frac23\sum_{i<j}(e_{xi}e_{xj}+e_{yi}e_{yj})\nonumber\\
 &+\frac23\sum_{\mathrm{cyc}(i,j,k)}e_{zk}
 (e_{xi}e_{xj}+e_{yi}e_{yj})\bigg].
 \label{eq:W-W}
\end{align}
Se obtiene $\Fstar^{(S)}=-2/9$ para los tres cortes. El mapa no modificado
presenta menor autovalor $-0.1547005$. $H_1$ permanece positiva, pero
$\lambda_{\min}(H_2)=-0.0411283$ certifica GME con momentos hasta quinto orden.
El SDP da $\Ngstar=0.4428090$.

\subsection{Extensión explícita a cuatro qubits}

Para demostrar que la construcción no depende de una peculiaridad de tres
qubits, traducimos al álgebra Wigner-star el ejemplo GHZ$_4$ de
Mukherjee \emph{et al.}~\cite{Mukherjee2025GMEMoments} y lo introducimos
directamente mediante
\begin{align}
 W_{\mathrm{GHZ}_4}=\frac1{16}\bigg[&1+\sum_{i<j}e_{zi}e_{zj}
 +e_{zA}e_{zB}e_{zC}e_{zD}+\prod_i e_{xi}+\prod_i e_{yi}\nonumber\\
 &-\sum_{i<j}e_{yi}e_{yj}\!\prod_{r\notin\{i,j\}}e_{xr}\bigg].
 \label{eq:W-GHZ4}
\end{align}
La fórmula contiene las cadenas $I/Z$ con un número par de factores $Z$ y las
cadenas $X/Y$ con un número par de factores $Y$, con signo
$(-1)^{N_Y/2}$. Las cuatro particiones $1|3$ y las tres particiones $2|2$ dan
$\Fstar^{(S)}=-1/4$ y negatividad $1/2$. El mapa modificado tiene menor
autovalor $-1/2$: $H_1$ es positiva, mientras
$\lambda_{\min}(H_2)=-0.6337602$ certifica GME. El SDP produce
$\Ngstar=0.499999993$.

\subsection{Contraejemplo biseparable NPT en todos los cortes}

Definimos directamente
\begin{equation}
 W_{\mathrm{bs}}
 =\frac1{24}\sum_{\mathrm{cyc}(i,j,k)}
 \left(1+e_{xi}e_{xj}-e_{yi}e_{yj}+e_{zi}e_{zj}\right)(1+e_{zk}).
 \label{eq:W-biseparable}
\end{equation}
Cada sumando es el símbolo de Bell $\Phi^+$ en un par y $|0\rangle$ en el
qubit restante; por construcción, la mezcla es biseparable. Sin embargo,
\begin{equation}
 \Fstar^{(A)}=\Fstar^{(B)}=\Fstar^{(C)}=-\frac1{54}<0.
\end{equation}
La capa de mapas no produce falso positivo: el menor autovalor de
$\mathcal M_T^{(3)}[W_{\mathrm{bs}}]$ es $2/3$, las matrices
$H_1,H_2,H_3$ son positivas y $\Ngstar=0$ dentro de la precisión del solver.
Este ejemplo demuestra por qué los cortes deben conservarse como diagnóstico
local y no como clasificador GME mixto.

\subsection{Ruido blanco}

Para $\psi\in\{\mathrm{GHZ},\mathrm W\}$ usamos la entrada explícita
\begin{equation}
 W_{\psi}(\mu)=\mu W_\psi+(1-\mu)\frac18,
 \qquad 0\leq\mu\leq1.
 \label{eq:white-noise-W}
\end{equation}
La Tabla~\ref{tab:examples} resume seis pruebas. Un guion en la columna Hankel
significa que ninguna de $H_1,H_2,H_3$ detecta, no que el estado sea
biseparable.

\begin{table}[ht]
\centering
\caption{Resultados obtenidos desde coeficientes explícitos de $W$. En estados
simétricos se reporta el valor común de todos los cortes no redundantes.}
\label{tab:examples}
\small
\begin{tabular}{lrrrrl}
\toprule
Entrada & $\Fstar$ & $\lambda_{\min}(\mathcal M W)$ & primer $H_\ell$ & $\Ngstar$ & conclusión \\
\midrule
GHZ$_3$ & $-0.250000$ & $-0.500000$ & $H_1$ & $0.500000$ & GME \\
W$_3$ & $-0.222222$ & $-0.154701$ & $H_2$ & $0.442809$ & GME \\
GHZ$_4$ & $-0.250000$ & $-0.500000$ & $H_2$ & $0.500000$ & GME \\
$W_{\mathrm{bs}}$ & $-0.018519$ & $+0.666667$ & -- & $0$ & biseparable \\
GHZ$_3(0.50)$ & $-0.008331$ & $+0.437500$ & -- & $0.062500$ & GME por SDP \\
W$_3(0.90)$ & $-0.158938$ & $-0.001730$ & -- & $0.345225$ & GME; Hankel $\ell\leq3$ inconcluso \\
\bottomrule
\end{tabular}
\end{table}

Los umbrales numéricos de visibilidad se muestran en la
Tabla~\ref{tab:thresholds}. El mapa completo detecta antes o igual que cualquier
truncación. Para W$_3$, $H_3$ se aproxima al umbral del mapa pero conserva una
pequeña brecha; esto cuantifica que un número finito de momentos puede ser
inconcluso incluso cuando el mapa completo ya es no positivo.

\begin{table}[ht]
\centering
\caption{Menor $\mu$ detectado para las familias de la
Ec.~\eqref{eq:white-noise-W}. ``No detecta'' significa que el nivel no detecta
ni siquiera en $\mu=1$.}
\label{tab:thresholds}
\small
\begin{tabular}{lcccc}
\toprule
Familia / mapa & mapa completo & $H_1$ & $H_2$ & $H_3$ \\
\midrule
GHZ$_3$ / modificado & $0.733333334$ & $0.934266784$ & $0.733333335$ & $0.733333334$ \\
W$_3$ / transposición & $0.898868744$ & no detecta & $0.953490622$ & $0.900474281$ \\
\bottomrule
\end{tabular}
\end{table}

\section{Validación numérica y reproducibilidad}

El archivo \path{scripts/verify_detector_jerarquico.py} implementa la tabla
local \eqref{eq:star-table}, la reflexión, las potencias estrella, Newton, los
mapas GME y el SDP. Las entradas GHZ$_3$, W$_3$, GHZ$_4$ y biseparable se
codifican mediante las Ecs.~\eqref{eq:W-GHZ}, \eqref{eq:W-W},
\eqref{eq:W-GHZ4} y \eqref{eq:W-biseparable}. El archivo
\path{scripts/validate_results.py} constituye una puerta separada que falla
ante inconsistencias o falsos positivos conocidos. Los comandos son
{\footnotesize
\begin{verbatim}
python -m pip install -r requirements.txt
python scripts/verify_detector_jerarquico.py
python scripts/validate_results.py
python scripts/verify_four_qubit_extension.py
\end{verbatim}
}
Las salidas completas se guardan en
{\small\path{results/detector_jerarquico_checks.json}} y
{\small\path{results/four_qubit_extension.json}}.\par
Las comprobaciones incluyen:
\begin{enumerate}
\item reconstrucción de una entrada física y retorno de coeficientes;
\item igualdad entre momentos estrella y trazas de potencias calculadas por una
vía matricial independiente;
\item igualdad entre el mapa aplicado a coeficientes SW y su acción matricial;
\item ausencia de falsos positivos sobre la mezcla biseparable;
\item evaluación de las siete biparticiones no redundantes de GHZ$_4$;
\item solución de los SDP con CLARABEL y respaldo previsto con SCS.
\end{enumerate}
La puerta de integridad pasa y el máximo error entre representaciones es
$9.094947\times10^{-13}$. Este error aparece en momentos de orden alto del mapa;
los momentos PT por corte concuerdan típicamente a nivel de redondeo
$10^{-16}$.

\section{Discusión}

\subsection{Qué mejora respecto al detector por determinante}

La mejora no consiste en reemplazar un determinante por otro. En la primera
capa, la familia completa $\{e_k^{(S)}\}$ evita los fallos de paridad y rango de
$\det(\Theta_SW)$. En la segunda, la condición
$\lambda_{\min}(H_\ell)<0$ usa toda la positividad de Hankel y evita repetir el
mismo fallo de paridad sobre $\det H_\ell$. En la tercera, el problema global se
formula sobre el casco convexo de estados PPT por cortes, que contiene la
información de descomposición ausente en una lista de espectros PT separados.

\subsection{Alcance de cada cantidad}

$\Fstar^{(S)}$ es un funcional espectral exacto de NPT por corte, no una medida
LOCC universal. $\mathcal G_{\star,\ell}^{(\mathcal M)}$ es un certificado suficiente y
depende de la elección del mapa y del orden; cero significa ``no detectado''.
$\Ngstar$ sí es un monótono heredado del SDP de negatividad genuina, pero su
positividad caracteriza estar fuera de las mezclas PPT, no todo el conjunto GME.
Esta separación de roles es deliberada.

\subsection{Limitaciones abiertas}

\begin{enumerate}
\item Los mapas de transposición usados aquí son descomponibles y no detectan
GME PPT. Mapas indecomponibles podrían ampliar la segunda capa, pero requieren
una construcción SW y una validación propia.
\item Una mezcla PPT puede ser GME; por ello $\Ngstar=0$ no es una prueba de
biseparabilidad salvo en clases especiales, como los estados
permutacionalmente invariantes de tres qubits.
\item La selección entre el mapa estándar y el modificado usa conocimiento de
familias GHZ/W. Una biblioteca de mapas o una optimización sobre mapas admisibles
es un problema posterior.
\item La adquisición de $W$ completa requiere $4^N-1$ parámetros reales. El
beneficio experimental de momentos truncados debe demostrarse con un protocolo
de medición concreto y análisis de varianza.
\item Newton puede ser numéricamente mal condicionado cerca de la frontera PPT
o para $N$ grande. Intervalos certificados, SDP de momentos o reconstrucción
espectral estable deben acompañar datos ruidosos.
\item Todo el formalismo numérico usa el kernel de la Ec.~\eqref{eq:kernel}. Un
cambio de kernel o medida exige transformar la normalización de $\Tfun_N$ y la
tabla estrella.
\end{enumerate}

\section{Conclusiones}

Hemos construido una versión multiqubit del detector Wigner-star que distingue
tres preguntas: NPT por corte, certificación GME con pocos momentos y exclusión
global del conjunto de mezclas PPT. La primera capa es exacta cuando se conocen
los $2^N$ momentos requeridos. La prueba GHZ$_4$ verifica explícitamente las
siete biparticiones no redundantes sin reconstruir primero $\rho$. La segunda
produce el funcional
$\mathcal G_{\star,\ell}$ a partir del menor autovalor de Hankel y puede certificar GHZ y
W con momentos de tercer y quinto orden, respectivamente. La tercera realiza el
SDP de testigos completamente descomponibles directamente sobre símbolos SW y
evita el falso positivo de la mezcla biseparable NPT en todos los cortes.

La formulación conserva como dato operacional a $W$ y hace explícita cada
normalización. Su principal valor es organizativo y metodológico: integra
espectroscopia PT, mapas GME y geometría convexa en una sola álgebra de fase sin
presentar como universal lo que solo es suficiente. El siguiente paso con mayor
valor científico es construir mapas indecomponibles en la base SW y comparar, con
un modelo de medición realista, el costo de estimar $\mathcal G_{\star,\ell}$ frente a
tomografía completa y testigos convencionales.

\section*{Disponibilidad de datos y código}

El código, los parámetros de los ejemplos y las salidas JSON están disponibles
en el repositorio público
\url{https://github.com/enrik03/wigner-star-multiqubit-entanglement}. La versión
asociada al envío se archivará en un repositorio de preservación para asignarle
un DOI. No se emplearon datos experimentales.

\section*{Declaraciones}

Los autores declaran no tener intereses en competencia. No se recibió
financiación específica para este estudio. E.E.C.B. y O.J.B.M. contribuyeron a
la conceptualización; E.E.C.B. desarrolló la implementación y la primera
redacción; ambos autores deberán revisar y aprobar la versión sometida.

{\small
\bibliographystyle{unsrt}
\bibliography{references}
}

\end{document}
